\documentclass[11pt]{article}
\usepackage[a4paper,margin=1in]{geometry}
\usepackage{amsmath,amssymb,amsthm,mathtools}
\usepackage{graphicx}
\usepackage{microtype}
\usepackage{enumitem}
\usepackage[hidelinks]{hyperref}

\newtheorem{theorem}{Theorem}
\newtheorem{lemma}{Lemma}
\newcommand{\F}{\mathbb F}
\newcommand{\C}{\mathbb C}
\newcommand{\Ran}{\operatorname{Ran}}
\newcommand{\Mcb}{M_{\mathrm{cb}}}

\title{Ordinary Norm-One Idempotent Multipliers of a Discrete Fourier Algebra Are Coset Indicators}
\author{Candidate full answer to the question in arXiv:0806.4643}
\date{11 August 2026}

\begin{document}
\maketitle

\begin{center}
\fbox{\parbox{0.92\textwidth}{\textbf{Status.} Candidate full proof, likely valid, pending human review. The source question has a negative answer, and the proof gives a classification for every discrete group.}}
\end{center}

\section{Source question and answer}

B. E. Forrest and V. Runde, \emph{Norm one idempotent cb-multipliers with
applications to the Fourier algebra in the cb-multiplier norm},
arXiv:0806.4643 \cite{source}, characterize norm-one idempotents in
$\Mcb A(G)$.  In the remark following Theorem~1.1 (PDF page~4), they note
that there are sets $E$ in the free group $\F_2$ for which
\[
 \chi_E\in MA(\F_2)\setminus \Mcb A(\F_2),
\]
and ask whether such an $E$ can be chosen with ordinary multiplier norm
$\|\chi_E\|_M=1$.

\begin{figure}[ht]
  \centering
  \includegraphics[width=\textwidth]{figures/open_problem_crop.png}
  \caption{The open question in the source PDF, page~4.}
\end{figure}

The answer is no.  The following result is stronger because it applies to
every discrete group.

\begin{theorem}\label{thm:main}
Let $G$ be a discrete group and let $\varnothing\ne E\subset G$. The following
are equivalent:
\begin{enumerate}[label=\textup{(\roman*)}]
 \item $\chi_E\in MA(G)$ and $\|\chi_E\|_M=1$;
 \item $E$ is a left coset of a subgroup of $G$;
 \item $\chi_E\in \Mcb A(G)$ and $\|\chi_E\|_{\Mcb}=1$.
\end{enumerate}
Consequently no set requested in the source question exists.
\end{theorem}

\section{Proof intuition}

Translate $E$ so that it contains the identity.  On the dual
$A(G)^*=VN(G)$, its multiplier becomes a norm-one projection which fixes the
identity.  That makes the projection positive.  The discrete group von
Neumann algebra has a faithful finite trace, and the Fourier-diagonal
projection preserves it.  Kadison's inequality then has no room to be strict
on the range: every self-adjoint range element has its square in the range.
Thus the range is closed under the ambient Jordan product.

The range is spanned weak-* by precisely the group unitaries indexed by the
translated set.  Jordan multiplying two such unitaries produces
$\lambda(st)+\lambda(ts)$.  Since group unitaries are linearly independent,
both products must still be supported.  Together with adjoint closure, this
turns the translated set into a subgroup.  Undoing the translation gives a
coset, whose indicator is visibly a coefficient of a quasi-regular unitary
representation.

\section{Operator-algebra lemmas}

We use the canonical faithful normal trace $\tau$ on $VN(G)$, determined by
$\tau(\lambda(s))=0$ for $s\ne e$ and $\tau(1)=1$.

\begin{lemma}\label{lem:positive}
If $T:\mathcal A\to\mathcal B$ is a unital linear contraction between unital
C*-algebras, then $T$ is positive.
\end{lemma}

\begin{proof}
For every state $\omega$ on $\mathcal B$, the functional $\omega\circ T$ is
unital and has norm at most one.  Its norm is in fact one, and a unital
norm-one functional on a C*-algebra is positive.  Thus
$\omega(T(a))\ge0$ whenever $a\ge0$.  States separate the positive cone of
$\mathcal B$, so $T(a)\ge0$.
\end{proof}

\begin{lemma}\label{lem:jordan}
Let $\mathcal M$ be a finite von Neumann algebra with faithful normal trace
$\tau$, and let $P:\mathcal M\to\mathcal M$ be a unital positive projection
such that $\tau\circ P=\tau$. Then $\Ran P$ is self-adjoint and is closed under
the ambient Jordan product
\[
 u\circ v=\frac{uv+vu}{2}.
\]
\end{lemma}

\begin{proof}
A positive linear map preserves adjoints, so $\Ran P$ is self-adjoint.  If
$a=a^*\in\Ran P$, Kadison's inequality for unital positive maps
\cite{kadison} gives
\[
 P(a^2)\ge P(a)^2=a^2.
\]
The positive difference $P(a^2)-a^2$ has trace zero because $P$ preserves
$\tau$.  Faithfulness therefore forces $P(a^2)=a^2$, so $a^2\in\Ran P$.
For self-adjoint $a,b\in\Ran P$,
\[
 ab+ba=(a+b)^2-a^2-b^2\in\Ran P.
\]
Every range element is a complex linear combination of two self-adjoint range
elements, and bilinearity finishes the proof.
\end{proof}

\section{Proof of the classification}

\begin{proof}[Proof of Theorem~\ref{thm:main}]
Assume (i). Choose $x\in E$ and put $H=x^{-1}E$, so $e\in H$. Left
translation of functions is an isometry of $A(G)$, and conjugating a
multiplication operator by translations shows that
\[
 \chi_H\in MA(G),\qquad \|\chi_H\|_M=1.
\]

Let
\[
 P=M_{\chi_H}^*:VN(G)\longrightarrow VN(G).
\]
For each $s\in G$, duality with point evaluation gives
\begin{equation}\label{eq:diagonal}
 P(\lambda(s))=\chi_H(s)\lambda(s).
\end{equation}
The map $P$ is normal, $P^2=P$, $\|P\|=1$, and $P(1)=1$ because $e\in H$.
Lemma~\ref{lem:positive} makes $P$ positive.  In particular $P$ preserves
adjoints.  If $s\in H$, then
\[
 \chi_H(s^{-1})\lambda(s^{-1})
 =P(\lambda(s)^*)=P(\lambda(s))^*=\lambda(s^{-1}),
\]
so $s^{-1}\in H$.

Equation~\eqref{eq:diagonal} also shows on the group algebra that
$\tau(P(a))=\tau(a)$, since $\chi_H(e)=1$.  Both functionals are normal and
the group algebra is weak-* dense in $VN(G)$, hence $\tau\circ P=\tau$ on all
of $VN(G)$. Lemma~\ref{lem:jordan} now applies.

Let $s,t\in H$. Then $\lambda(s),\lambda(t)\in\Ran P$, and Jordan closure
gives
\[
 z:=\lambda(st)+\lambda(ts)\in\Ran P.
\]
Therefore $P(z)=z$. If $st=ts$, equation~\eqref{eq:diagonal} immediately
gives $\chi_H(st)=1$. If $st\ne ts$, it gives
\[
 \chi_H(st)\lambda(st)+\chi_H(ts)\lambda(ts)
 =\lambda(st)+\lambda(ts).
\]
Distinct group unitaries are linearly independent (apply them to
$\delta_e\in\ell^2(G)$), so again $\chi_H(st)=1$; indeed both coefficients
are one. Thus $H$ contains $e$ and is closed under products and inverses.
It is a subgroup, and $E=xH$. This proves (i)$\Rightarrow$(ii).

For (ii)$\Rightarrow$(iii), let $E=xH$ and let $\pi$ be the quasi-regular
representation on $\ell^2(G/H)$. Then
\[
 \chi_{xH}(s)=\langle \pi(s)\delta_H,\delta_{xH}\rangle.
\]
Hence $\chi_E\in B(G)$ with $B(G)$-norm at most one. Since it takes the value
one, the norm is exactly one. The canonical contractive inclusion
$B(G)\hookrightarrow\Mcb A(G)$ gives cb multiplier norm at most one, and a
nonzero idempotent operator has norm at least one. Thus the cb norm is one.

Finally, the contractive inclusion $\Mcb A(G)\hookrightarrow MA(G)$ gives
(iii)$\Rightarrow$(i), again using that a nonzero idempotent has norm at least
one.
\end{proof}

\section{Verification, scope, and novelty}

The proof uses ordinary multiplier contractivity only; complete boundedness
enters after the support has already been classified as a coset.  Discreteness
is used exactly in the availability of the canonical faithful finite trace on
$VN(G)$.  Therefore this packet does not assert the analogous classification
for nondiscrete locally compact groups, and it does not address the distinct
amenability questions later in the source paper.

As a finite check, the exact abelian Fourier formula was used to enumerate all
8,177 nonempty subsets of cyclic groups $C_n$, $2\le n\le12$.  Ordinary
multiplier norm one occurred exactly for cosets.  The code is included in
\texttt{code/check\_cyclic\_groups.py}; the proof does not depend on it.

The run's lightweight indexes were searched for arXiv id 0806.4643 and the
core ordinary-multiplier/coset terms.  A bounded external search on 11 August
2026 checked the exact source wording, close norm-one idempotent variants,
the source's citation graph, and metadata for the closest later work on
idempotent multiplier norms and contractive projections.  No exact answer or
statement of the all-discrete-group classification was found.  This was not
an exhaustive literature review; originality is therefore a candidate claim
pending expert checking.

The recommended human-review focus is the short trace-to-Jordan argument and
the Fourier-diagonal identification \eqref{eq:diagonal}.  No conditional
dependency remains.

\enlargethispage{3\baselineskip}
{\small
\begin{thebibliography}{9}
\bibitem{source}
B. E. Forrest and V. Runde,
\newblock Norm one idempotent cb-multipliers with applications to the Fourier
algebra in the cb-multiplier norm,
\newblock arXiv:0806.4643 (2008),
\href{https://arxiv.org/abs/0806.4643}{arXiv:0806.4643}.

\bibitem{kadison}
R. V. Kadison,
\newblock A generalized Schwarz inequality and algebraic invariants for
operator algebras,
\newblock \emph{Ann. of Math.} (2) \textbf{56} (1952), 494--503.
\end{thebibliography}
}

\end{document}
